\chapter{Classification périodique des éléments}

\noindent\textit{Plus d'une centaine d'éléments chimiques différents existent dans la
nature. Pour s'y retrouver, les scientifiques les ont rangés dans un grand tableau — la
classification périodique — qui permet de prévoir leurs propriétés d'un simple coup d'œil.}
\vspace{8pt}

\connecteBanner

\section{Élément chimique et numéro atomique}

\begin{definition}[Élément chimique]
Un \textbf{élément chimique} regroupe tous les atomes (ou ions dérivés de ces atomes) qui
possèdent le même nombre de protons dans leur noyau. Ce nombre de protons est appelé
\textbf{numéro atomique}, noté $Z$.
\end{definition}

\begin{propriete}[Numéro atomique et électrons]
Pour un atome \textbf{neutre}, le numéro atomique $Z$ est aussi égal au nombre d'électrons
(neutralité électrique, chapitre précédent). Chaque élément possède un numéro atomique
\textbf{unique}, qui le distingue de tous les autres.
\end{propriete}

\begin{illus}
\begin{tikzpicture}[scale=0.72,every node/.style={font=\scriptsize}]
\foreach \x/\s/\z in {0/H/1,1/He/2,2/Li/3,3/Be/4,4/B/5,5/C/6,6/N/7,7/O/8,8/F/9,9/Ne/10}{
  \draw[onbo,line width=1pt,fill=onbsoft] (\x*1.05,0) rectangle ++(0.9,0.9);
  \node at (\x*1.05+0.45,0.55) {\s};
  \node[onbmuted] at (\x*1.05+0.45,0.2) {\z};
}
\draw[->,onbmuted,line width=1pt] (0,-0.5) -- (9.9,-0.5);
\node[onbmuted,below] at (4.9,-0.7) {\small numéro atomique $Z$ croissant};
\end{tikzpicture}
\fig{Figure — Les dix premiers éléments, rangés par numéro atomique croissant.}
\end{illus}

\begin{exemple}
Le sodium a pour numéro atomique $Z=11$ : un atome de sodium possède $11$ protons et, s'il
est neutre, $11$ électrons.
\end{exemple}

\section{La classification périodique}

\begin{propriete}[Règle de la classification périodique]
Les éléments chimiques sont rangés dans un tableau appelé \textbf{classification
périodique}, par numéro atomique \textbf{croissant}, de gauche à droite puis de haut en
bas. Les éléments sont organisés en \textbf{lignes} (appelées périodes) et en
\textbf{colonnes} (appelées familles ou groupes).
\end{propriete}

\begin{propriete}[Éléments d'une même colonne]
Les éléments d'une même \textbf{colonne} appartiennent à la même \textbf{famille
chimique} : ils ont des propriétés chimiques \textbf{semblables}.
\end{propriete}

\begin{illus}
\begin{tikzpicture}[scale=0.62,every node/.style={font=\scriptsize}]
% ligne 1
\draw[onbgreen,line width=1pt,fill=onbgreen!15] (0,1.8) rectangle ++(0.9,0.9);
\node at (0.45,2.35) {H};
\node[onbmuted] at (0.45,2.0) {1};
\draw[onbblue,line width=1pt,fill=onbblue!15] (8.5,1.8) rectangle ++(0.9,0.9);
\node at (8.95,2.35) {He};
\node[onbmuted] at (8.95,2.0) {2};
% ligne 2
\draw[onbgreen,line width=1pt,fill=onbgreen!15] (0,0.8) rectangle ++(0.9,0.9);
\node at (0.45,1.35) {Li};
\node[onbmuted] at (0.45,1.0) {3};
\foreach \x/\s/\z/\i in {1/Be/4/1,2/B/5/2,3/C/6/3,4/N/7/4,5/O/8/5}{
  \draw[onbline,line width=0.8pt,fill=white] (\x*1.05,0.8) rectangle ++(0.9,0.9);
  \node at (\x*1.05+0.45,1.35) {\s};
  \node[onbmuted] at (\x*1.05+0.45,1.0) {\z};
}
\draw[onbo,line width=1pt,fill=onbo!15] (6.3,0.8) rectangle ++(0.9,0.9);
\node at (6.75,1.35) {F};
\node[onbmuted] at (6.75,1.0) {9};
\draw[onbblue,line width=1pt,fill=onbblue!15] (8.5,0.8) rectangle ++(0.9,0.9);
\node at (8.95,1.35) {Ne};
\node[onbmuted] at (8.95,1.0) {10};
% ligne 3
\draw[onbgreen,line width=1pt,fill=onbgreen!15] (0,-0.2) rectangle ++(0.9,0.9);
\node at (0.45,0.35) {Na};
\node[onbmuted] at (0.45,0.0) {11};
\foreach \x/\s/\z/\i in {1/Mg/12/1,2/Al/13/2,3/Si/14/3,4/P/15/4,5/S/16/5}{
  \draw[onbline,line width=0.8pt,fill=white] (\x*1.05,-0.2) rectangle ++(0.9,0.9);
  \node at (\x*1.05+0.45,0.35) {\s};
  \node[onbmuted] at (\x*1.05+0.45,0.0) {\z};
}
\draw[onbo,line width=1pt,fill=onbo!15] (6.3,-0.2) rectangle ++(0.9,0.9);
\node at (6.75,0.35) {Cl};
\node[onbmuted] at (6.75,0.0) {17};
\draw[onbblue,line width=1pt,fill=onbblue!15] (8.5,-0.2) rectangle ++(0.9,0.9);
\node at (8.95,0.35) {Ar};
\node[onbmuted] at (8.95,0.0) {18};
% legend
\node[onbgreen,left] at (-0.3,-0.9) {\scriptsize\textbullet\ alcalins};
\node[onbo] at (3.4,-0.9) {\scriptsize\textbullet\ halogènes};
\node[onbblue,right] at (6.3,-0.9) {\scriptsize\textbullet\ gaz nobles};
\end{tikzpicture}
\fig{Figure — Extrait du tableau périodique (3 premières périodes) : trois familles
chimiques mises en évidence.}
\end{illus}

\begin{exemple}
Le fluor $\mathrm{F}$ ($Z=9$) et le chlore $\mathrm{Cl}$ ($Z=17$) sont dans la même colonne
: ils appartiennent tous deux à la famille des \textbf{halogènes} et ont des propriétés
chimiques semblables (ce sont des non-métaux très réactifs).
\end{exemple}

\begin{remarque}
Les \textbf{gaz nobles} (hélium, néon, argon...) forment la famille la plus à droite du
tableau : ce sont des gaz très peu réactifs, presque inertes chimiquement.
\end{remarque}

\section{Utiliser la classification périodique}

\begin{propriete}[Utilité pratique]
Connaître la position d'un élément dans la classification permet de \textbf{prévoir} ses
propriétés chimiques sans avoir à les mesurer, simplement en les comparant à celles des
autres éléments de sa colonne (même famille).
\end{propriete}

\begin{exemple}
Le sodium $\mathrm{Na}$ ($Z=11$) et le potassium $\mathrm{K}$ ($Z=19$) sont tous deux dans
la colonne des \textbf{alcalins} : comme le sodium, le potassium est un métal mou, très
réactif, qui réagit vivement avec l'eau.
\end{exemple}

\begin{remarque}
Le numéro atomique augmente toujours de $1$ en $1$ en parcourant une période de gauche à
droite : entre deux éléments \textbf{voisins} d'une même ligne, il n'y a jamais d'élément
manquant.
\end{remarque}

% ═══════════════════════════════════════════════════════════════
\section{L'essentiel à retenir}

\begin{synthese}
\textbf{Élément chimique.} Ensemble des atomes ayant le même numéro atomique $Z$ (nombre de
protons).\\[4pt]
\textbf{Numéro atomique.} Pour un atome neutre : $Z=$ nombre de protons $=$ nombre
d'électrons.\\[4pt]
\textbf{Classification périodique.} Éléments rangés par $Z$ croissant, en lignes
(périodes) et colonnes (familles).\\[4pt]
\textbf{Famille chimique.} Éléments d'une même colonne $\Rightarrow$ propriétés chimiques
semblables (ex. alcalins, halogènes, gaz nobles).\\[4pt]
\textbf{Piège fréquent.} Deux éléments d'une même \textbf{ligne} n'ont pas forcément de
propriétés semblables ; c'est la \textbf{colonne} qui détermine la famille chimique.
\end{synthese}

% ═══════════════════════════════════════════════════════════════
\section{Méthodes \& savoir-faire}

\methode{Déterminer le nombre d'électrons d'un atome neutre à partir de son numéro
atomique}
Pour un atome neutre, le nombre d'électrons est toujours égal au numéro atomique $Z$
(nombre de protons).
\begin{exemple}
Le calcium a pour numéro atomique $Z=20$ : un atome de calcium neutre possède $20$
électrons.
\end{exemple}

\methode{Comparer deux éléments à l'aide de leur numéro atomique}
Le numéro atomique le plus petit correspond à l'élément situé le plus à gauche et/ou le
plus haut dans la classification.
\begin{exemple}
Le carbone ($Z=6$) est situé avant l'oxygène ($Z=8$) dans la classification.
\end{exemple}

\methode{Identifier la famille chimique d'un élément}
Repérer la colonne de l'élément dans la classification : tous les éléments de cette
colonne forment la même famille chimique.
\begin{exemple}
Le chlore ($Z=17$) est dans la colonne des halogènes, comme le fluor ($Z=9$).
\end{exemple}

\methode{Prévoir une propriété par comparaison au sein d'une famille}
Si deux éléments sont dans la même colonne, ils partagent des propriétés chimiques
similaires : on peut prévoir le comportement de l'un en connaissant celui de l'autre.
\begin{exemple}
Le sodium réagit vivement avec l'eau ; le potassium, autre alcalin, réagit lui aussi
vivement (et même plus violemment) avec l'eau.
\end{exemple}

\methode{Reconnaître un gaz noble}
Les gaz nobles occupent la dernière colonne (la plus à droite) de la classification et se
caractérisent par une très faible réactivité chimique.
\begin{exemple}
L'hélium ($Z=2$), le néon ($Z=10$) et l'argon ($Z=18$) sont des gaz nobles : ils
n'entrent quasiment dans aucune réaction chimique.
\end{exemple}

% ═══════════════════════════════════════════════════════════════
\section{Exercices d'application}
\begin{multicols}{2}

\rubrique{Numéro atomique}
\exo{}\diff{1} Le sodium a pour numéro atomique $Z=11$. Combien de protons et combien
d'électrons possède un atome de sodium neutre ?

\exo{}\diff{1} Le fer a pour numéro atomique $Z=26$. Combien d'électrons possède un atome
de fer neutre ?

\exo{}\diff{2} Qu'est-ce qu'un élément chimique ?

\exo{}\diff{2} Deux atomes ont respectivement $8$ et $9$ protons dans leur noyau.
Appartiennent-ils au même élément chimique ? Justifier.

\rubrique{Structure de la classification}
\exo{}\diff{1} Comment les éléments sont-ils rangés dans la classification périodique ?

\exo{}\diff{2} Le carbone a pour numéro atomique $6$ et l'azote $7$. Lequel des deux est
placé avant l'autre dans la classification ?

\exo{}\diff{2} Que représente une colonne de la classification périodique ?

\exo{}\diff{1} Que représente une ligne (période) de la classification périodique ?

\rubrique{Familles chimiques}
\exo{}\diff{1} Citer une famille chimique de la classification périodique.

\exo{}\diff{2} Le fluor et le chlore sont dans la même colonne. Que peut-on en déduire sur
leurs propriétés chimiques ?

\exo{}\diff{2} Pourquoi les gaz nobles sont-ils peu utilisés dans les réactions
chimiques ?

\exo{}\diff{2} L'hélium et le néon appartiennent-ils à la même famille chimique ?
Justifier.

\rubrique{Utilisation du tableau}
\exo{}\diff{2} Le sodium réagit vivement avec l'eau. Que peut-on prévoir pour le
potassium, qui appartient à la même famille ?

\exo{}\diff{2} Un élément a pour numéro atomique $Z=17$. À quelle famille appartient-il ?

\exo{}\diff{3} Deux éléments appartiennent à la même ligne de la classification. Ont-ils
nécessairement des propriétés chimiques semblables ? Justifier.

\exo{}\diff{2} Expliquer en une phrase l'intérêt pratique de la classification
périodique.

% ═══════════════════════════════════════════════════════════════
\end{multicols}

\section{Exercices d'entraînement}
\begin{multicols}{2}

\rubrique{Numéro atomique}
\exo{}\diff{1} Le magnésium a pour numéro atomique $Z=12$. Combien de protons et
combien d'électrons possède un atome de magnésium neutre ?

\exo{}\diff{1} L'aluminium a pour numéro atomique $Z=13$. Combien d'électrons possède un
atome d'aluminium neutre ?

\exo{}\diff{2} Un atome possède $16$ électrons et est électriquement neutre. Quel est son
numéro atomique ?

\exo{}\diff{2} Le potassium a pour numéro atomique $19$ et le calcium $20$. Lequel des
deux atomes possède le plus d'électrons ?

\exo{}\diff{3} Deux atomes appartiennent au même élément chimique. Que peut-on affirmer
sur leur nombre de protons ?

\rubrique{Structure de la classification}
\exo{}\diff{1} La classification périodique range-t-elle les éléments par numéro
atomique croissant ou décroissant ?

\exo{}\diff{2} Le silicium a pour numéro atomique $14$ et le phosphore $15$. Lequel des
deux se trouve le plus à gauche dans sa ligne ?

\exo{}\diff{2} Comment appelle-t-on une colonne de la classification périodique ?

\exo{}\diff{1} Comment appelle-t-on une ligne de la classification périodique ?

\exo{}\diff{2} Vrai ou Faux : « Entre deux éléments voisins d'une même ligne, il peut
manquer un élément intermédiaire. »

\rubrique{Familles chimiques}
\exo{}\diff{1} Citer deux éléments de la famille des halogènes.

\exo{}\diff{2} Citer deux éléments de la famille des gaz nobles.

\exo{}\diff{2} Le lithium et le sodium sont dans la même colonne. À quelle famille
appartiennent-ils ?

\exo{}\diff{2} Pourquoi dit-on que les alcalins sont des métaux très réactifs ?

\exo{}\diff{3} Le brome appartient à la même famille que le chlore. Sachant que le chlore
est un non-métal réactif, que peut-on prévoir pour le brome ?

\rubrique{Utilisation du tableau}
\exo{}\diff{2} Un élément a pour numéro atomique $Z=10$. À quelle famille appartient-il
? Que peut-on en déduire sur sa réactivité ?

\exo{}\diff{2} Le potassium réagit très vivement avec l'eau. Que peut-on prévoir pour le
lithium, autre alcalin ?

\exo{}\diff{2} Expliquer pourquoi il est utile de connaître la famille chimique d'un
élément avant de l'utiliser dans une expérience.

\exo{}\diff{3} Un élève affirme : « Le carbone et l'azote, voisins dans le tableau, ont
forcément les mêmes propriétés chimiques. » Cette affirmation est-elle correcte ?
Justifier.

\exo{}\diff{2} Citer un élément de la famille des alcalins et un élément de la famille
des halogènes.

% ═══════════════════════════════════════════════════════════════
\end{multicols}

\section{Sujets type BEPC}

\sujetbac[Numéro atomique et familles chimiques]
\begin{enumerate}[label=\arabic*)]
\item Définis : élément chimique.
\item Le zinc a pour numéro atomique $Z=30$. Combien d'électrons possède un atome de zinc
neutre ?
\item Cite un élément de la famille des gaz nobles.
\item Le fluor et le chlore sont-ils dans la même famille chimique ? Justifier.
\end{enumerate}

\sujetbac[Lecture de la classification périodique]
\begin{enumerate}[label=\arabic*)]
\item Comment les éléments sont-ils rangés dans la classification périodique ?
\item Que représente une colonne de la classification ?
\item Le sodium ($Z=11$) et le potassium ($Z=19$) sont dans la même colonne. Que peut-on
prévoir sur leurs propriétés chimiques ?
\item Réponds par Vrai ou Faux : « Les gaz nobles sont très réactifs chimiquement. »
\end{enumerate}

\sujetbac[Utiliser la classification pour prévoir des propriétés]
\begin{enumerate}[label=\arabic*)]
\item Qu'est-ce que le numéro atomique d'un élément ?
\item Un atome possède $17$ électrons et est neutre. Quel est son numéro atomique ? À
quelle famille appartient-il ?
\item Le chlore réagit facilement avec les métaux. Que peut-on prévoir pour le brome,
autre halogène ?
\item Cite l'intérêt pratique de la classification périodique pour un chimiste.
\end{enumerate}

% ═══════════════════════════════════════════════════════════════
\section{Corrigés}
\begin{multicols}{2}

\corrige{de l'exercice 1}
$11$ protons et $11$ électrons.

\corrige{de l'exercice 2}
$26$ électrons.

\corrige{de l'exercice 3}
Un élément chimique regroupe tous les atomes ayant le même numéro atomique (même nombre
de protons).

\corrige{de l'exercice 4}
Non : ils ont un nombre de protons différent ($8$ et $9$), donc ils appartiennent à des
éléments chimiques différents.

\corrige{de l'exercice 5}
Par numéro atomique croissant, en lignes (périodes) et colonnes (familles).

\corrige{de l'exercice 6}
Le carbone ($Z=6$) est placé avant l'azote ($Z=7$).

\corrige{de l'exercice 7}
Une colonne représente une famille chimique : des éléments aux propriétés chimiques
semblables.

\corrige{de l'exercice 8}
Une ligne (période) regroupe des éléments consécutifs par numéro atomique croissant.

\corrige{de l'exercice 9}
Par exemple : la famille des halogènes (ou des alcalins, ou des gaz nobles).

\corrige{de l'exercice 10}
Ils ont des propriétés chimiques semblables : ce sont tous deux des non-métaux réactifs
(halogènes).

\corrige{de l'exercice 11}
Parce qu'ils sont chimiquement très peu réactifs (presque inertes).

\corrige{de l'exercice 12}
Oui : ils sont dans la même colonne (famille des gaz nobles).

\corrige{de l'exercice 13}
Le potassium devrait lui aussi réagir vivement avec l'eau (même famille que le sodium).

\corrige{de l'exercice 14}
Il appartient à la famille des halogènes.

\corrige{de l'exercice 15}
Non : ce sont des éléments voisins d'une même ligne, mais chaque colonne correspond à une
famille différente ; leurs propriétés ne sont pas nécessairement semblables.

\corrige{de l'exercice 16}
Elle permet de prévoir les propriétés chimiques d'un élément par comparaison avec les
autres éléments de sa famille.

\corrige{du sujet 1}
1) Un élément chimique regroupe tous les atomes ayant le même numéro atomique.\\
2) $30$ électrons.\\
3) Par exemple l'hélium, le néon ou l'argon.\\
4) Oui : ils sont dans la même colonne (famille des halogènes).

\corrige{du sujet 2}
1) Par numéro atomique croissant, en lignes et colonnes.\\
2) Une famille chimique : des éléments aux propriétés semblables.\\
3) Des propriétés chimiques semblables (tous deux alcalins, métaux réactifs).\\
4) Faux : les gaz nobles sont très peu réactifs.

\corrige{du sujet 3}
1) C'est le nombre de protons de son noyau (égal au nombre d'électrons pour un atome
neutre).\\
2) Numéro atomique $17$ ; famille des halogènes.\\
3) Le brome devrait lui aussi réagir facilement avec les métaux (même famille que le
chlore).\\
4) Elle permet de prévoir les propriétés d'un élément sans avoir à toutes les mesurer.

\end{multicols}
\exercicesBanner
